\documentstyle[%


righttag%


,11pt%


,amssymb%


,russian%


,izvru%


]{amsart}





\newcommand{\tl}[3]{\medskip\noindent{\bf#1}\ #2 \dotfill #3 \par} 


\pagestyle{empty}


\begin{document}


%\tableofcontents


\par


\vskip 2cm


\par


\bigskip





\bigskip


\bigskip


\par


\centerline{\Large СОДЕРЖАНИЕ}


\bigskip


\bigskip


%%%%%%%%%% Use \newline %%%%%%%%%%%











\tl{Агачев Ю.Р., Леонов А.И.}{Решение одного класса


интегродифференциальных\newline уравнений


методом механических квадратур}{3}


\tl{Воронецкая М.А., Иванов А.Г.}


{Почти периодическая задача Больца}{8}


\tl{Коняев Ю.А., Панфилов Н.Г.}


{Асимптотический анализ линейных периодических\newline систем 


одронодных дифференциальных уравнений


при наличии большого или малого\newline параметра}{25}


\tl{Корешков Н.А.}{Теорема Энгеля для алгебр лиевского типа}{30}


\tl{Кочкарев Б.С.} 


{Структурные свойства одного класса максимальных шпернеровых \newline


семейств подмножеств}{37}


\tl{Масальцев Л.А.}


{Бикасательное преобразование Бианки подмногообразия постоянной\newline


отрицательной кривизны $H^n$ евклидова пространства $R^{2n}$}{43}


\tl{Павленко В.Н., Чиж Е.А.}


{Сильно резонансные эллиптические вариационные


неравен-\newline ства с разрывными нелинейностями}{49}


\tl{Сильченко Ю.Т.}


{Полугруппы с неплотно заданным производящим оператором}{57}


\tl{Тимербаев М.Р.}


{Весовые оценки решения анизотропно вырождающегося уравнения\newline


с граничными условиями Неймана в точках вырождения}{63}











\bigskip


\centerline{\Large КРАТКИЕ~СООБЩЕНИЯ}


\bigskip





\tl{Ломакин Д.Е.}


{Оператор преобразования Радона целых функций многих 


комплексных\newline переменных}{77}








\vfill


\noindent\raisebox{2pt}{\copyright} Казанский госудаpственный унивеpситет


\end{document}


